%##########################################################################
%                                                                     
%                         LaTeX-Vorlage
%                                 
%            für Bachelor-, Studien-, Diplom- und Masterarbeiten
%                                                                  
%                      Ralf Kuebler 1996     
%                      
%                      ueberarbeitet Joerg Fehr 2009
%                      ueberarbeitet Christian Fischer 2011
%                      ergaenzt Thomas Kurz und Nicolai Wengert 2011
%					   ueberarbeitet Fedor Vlasov 2013 - 2016
%                      
%
%##########################################################################


%##########################################################################
\RequirePackage[l2tabu, orthodox]{nag} % warnings about deprecated packages
\documentclass[12pt,a4paper,twoside]{report}
\usepackage{etex}		% optimize memory usage
\usepackage{ucs}			% extended unicode character set
\usepackage[utf8x,utf8]{inputenc}
\usepackage{ngerman,amsmath,amssymb,amsthm}
\usepackage[ngerman]{babel}
\usepackage{cite} % for citing
\usepackage{acronym}	 % for list of acronyms
\usepackage{graphicx, epstopdf}	% adding pictures
\usepackage{pstricks-add}	% psticks graphics

\usepackage[utf8]{inputenc}
\usepackage[T1]{fontenc}
\usepackage{lmodern}
\usepackage{hyperref}
\usepackage{listings}
\usepackage{float}

\lstnewenvironment{code}[1][]%
{
   \noindent
   \minipage{\linewidth} 
   \vspace{0.5\baselineskip}
   \lstset{basicstyle=\ttfamily\footnotesize,frame=single,#1,breaklines=true}}
{\endminipage}



% This will fix the spacing problem in the list of algorithms
\usepackage{etoolbox}
\makeatletter
\patchcmd{\@chapter}% <cmd>
  {\chaptermark{#1}}% <search
  {\chaptermark{#1}%
   \addtocontents{loa}{\protect\addvspace{10\p@}}}% replace
  {}{}% <success><failure>
\makeatother

\usepackage[chapter]{algorithm}		% writing pseudo code
\usepackage[noend]{algpseudocode}	% writing pseudo code
\usepackage{subfigure}	% multiple figures with one caption
\usepackage{multirow} 	% multirows in tables
\usepackage{bm}
\usepackage{pbox}     	% e.g. for vertical text in tables
\usepackage{calc}
\usepackage{pdfpages} 	% insert pdf pages
\usepackage{nicefrac} 	% use nicely formated fractions in tex
\usepackage{listings} 	% add code snippets
\usepackage{hyperref}	% More descriptive referencing
\usepackage{tabularx}   % for more suffisticated table control}
\usepackage{booktabs}	% for scientific tables - no vertical lines
\usepackage{siunitx}		% S.I. units
\usepackage{caption}		% More control over caption position

\AtBeginDocument{\renewcommand{\chaptername}{}}% muss nach babel stehen! 
\usepackage[toctitles]{titlesec} 
\titleformat{\chapter}{\bfseries\huge}{\thechapter.\quad}{0em}{}

% This is a very usefull package from ITM (Uni Stuttgart)
% use this when you have to use vectors and matrices alot
% The file is located inside the template directory
\usepackage{mechb_symbol}

\pagestyle{headings}

% spacing between paragraphs
\setlength{\parskip}{1.5ex}

% spacing between captions and objects
\setlength\abovecaptionskip{5pt}
\setlength\belowcaptionskip{5pt}

% do not intend the first line of a paragraph
\setlength{\parindent}{0pt}

% sloppier spacing between words
\sloppy

% hopefully pictures are where you want them to be
\setcounter{topnumber}{20}\setcounter{bottomnumber}{20}
\setcounter{totalnumber}{20}
\renewcommand{\topfraction}{.9999}
\renewcommand{\bottomfraction}{.9999}
\renewcommand{\textfraction}{0}

%##########################################################################
% Put all your pictures in this folder:
\def \mygraphicspath {pictures/}
\graphicspath{{\mygraphicspath}} % Set the default folder for images

% some predefined macros for convinience
% ps/eps-file with caption:
\newcommand{\mypicture}[4]{
\begin{figure}[hptb]
\centering
\includegraphics[scale=1,angle=#4]{#1}
\caption{#2}\label{#3}
\end{figure}
}%newcommand%
% call: \mypicture{file}{caption}{label}{angle}

% ps/eps-file with caption and width
\newcommand{\mywidthpicture}[5]{
\begin{figure}[hptb]
\centering
\includegraphics[width=#5,angle=#4]{#1}
\caption{#2}\label{#3}
\end{figure}
}%newcommand%
% call: \mywidthpicture{file}{caption}{label}{angle}{width}

% ps/eps-file with caption and width
\newcommand{\mydennis}[5]{
\begin{figure}[hptb]
\centering
\includegraphics[width=#5,angle=#4]{#1}
\caption{#2}\label{#3}
\end{figure}
}%newcommand%
% call: \mywidthpicture{file}{caption}{label}{angle}{width}

% ps/eps-file with caption and scaling
\newcommand{\myscalepicture}[5]{
\begin{figure}[hptb]
\centering
\includegraphics[scale=#5,angle=#4]{#1}
\caption{#2}\label{#3}
\end{figure}
}%newcommand%
% call: \myscalepicture{file}{caption}{label}{angle}{scale}

% ps/eps-file with caption and scaling and float layout
\newcommand{\myfloatpicture}[5]{
\begin{figure}[H]
\centering
\includegraphics[scale=#5,angle=#4]{#1}
\caption{#2}\label{#3}
\end{figure}
}%newcommand%
% call: \myscalepicture{file}{caption}{label}{angle}{scale}

% This one is special: In incscape you have the possibility to export an eps-file
% with an .tex file containing the text of the image. You than get two files, e.g.:
% 1. pic.eps
% 2. pic.eps_tex
% Put both files in the ./picture folder. 
% For this command use the "pic.eps_tex" file
\newcommand{\myinkscapepicture}[4]{
\begin{figure}[hptb] 
\centering
\def\svgwidth{#4} 
\input{\mygraphicspath#1} 
\caption{#2} 
\label{#3}
\end{figure}
}
%call: \myinkscapepicture{file.eps_tex}{caption}{label}{width}


%##########################################################################
% math mode
\DeclareMathOperator\sign{sign} % sign function

% plus/minus symbol where one of the signs is in brackets
\newcommand*\myPM{\ensuremath{\substack{+\\[-0.25em]-}\,}}
\newcommand*\myMP{\ensuremath{\substack{-\\[-0.25em]+}\,}}
\newcommand*\myPMbM{\ensuremath{\substack{+\\[-0.25em]\mathsmaller(-\mathsmaller)}\,}}
\newcommand*\myPMbP{\ensuremath{\substack{\mathsmaller(+\mathsmaller)\\[-0.25em]-}\,}}


%###########################################################################
% insert code snippets
\definecolor{dkgreen}{rgb}{0,0.6,0}
\definecolor{gray}{rgb}{0.5,0.5,0.5}
\definecolor{mauve}{rgb}{0.58,0,0.82}

% code formatting
\lstset{basicstyle=\ttfamily,
  showstringspaces=false,
  commentstyle=\color{dkgreen},
  keywordstyle=\color{blue},
  stringstyle=\color{mauve},
  language=bash,
  xleftmargin=0pt,
  numbers=left,                 
  numbersep=5pt,
  numberstyle=\small\color{gray},
  frame=shadowbox,
  rulesepcolor=\color{gray},
  breakindent=20pt,
  belowcaptionskip=2\smallskipamount,
  breaklines
}

% Setup hyperref
\hypersetup{
  colorlinks   	= true, %Colors links instead of ugly boxes
  urlcolor     	= black, %Color for external hyperlinks
  linkcolor    	= black, %Color of internal links
  citecolor   	= black %Color of citations
}


% Namen von Referenzen anpassen
\addto\extrasngerman{\def\figureautorefname{Abbildung}}
\addto\extrasngerman{\def\tableautorefname{Tabelle}}

% Prevent footenotes from splitting
\interfootnotelinepenalty=10000

\begin{document}

%###########################################################################
%
%   Titelseite
%
%###########################################################################
\begin{titlepage}
%\vspace*{5mm}
\begin{center}
  %\vspace{10mm} 
         {\large  Projektarbeit im SoSe2017\\}
  \vspace{10mm}
         {\Large
	\textbf{Konzeption und Aufbau eines Handheld zum Auslesen des BCCH umliegender GSM Basisstationen }\\}
  \vspace{10mm}
         {\large von}\\
         {\large Dennis Dette [Mat.Nr.: 45311] 
          \\ Christian Kobiela [Matr.Nr.: 47571]}\\
  \vfill
  
  \vspace{5mm}
                  \includegraphics[width=5cm]{HsKaLogo.png}\\
  \vspace{5mm}
  \vspace{10mm}
         {\large 23.10.2017}
\end{center}
\end{titlepage}

\clearpage
\thispagestyle{empty}
\cleardoublepage

% content starts at the right page
\thispagestyle{empty}\cleardoublepage
% settings for 2-side printing
\evensidemargin=2pt
\oddsidemargin=40pt

\renewcommand{\baselinestretch}{1}\normalsize
\pagenumbering{roman}
%###########################################################################
%
%   Inhaltsverzeichnis
%
%   (wird automatisch erstellt; dieser File ist nicht zu ändern)
%
%###########################################################################
\tableofcontents

\chapter*{Abkürzungsverzeichnis}
\addcontentsline{toc}{chapter}{Abkürzungsverzeichnis}

\begin{acronym}[CSMA/CD]
	\acro{UML}{Unified Modelling Language}
	\acro{GSM}{Global System for Mobile Communications}
	\acro{BTS}{Base transceiver station}
	\acro{BCCH}{Broadcast control channel}
\end{acronym}
\addcontentsline{toc}{chapter}{Abbildungsverzeichnis}
\listoffigures
\addcontentsline{toc}{chapter}{Tabellenverzeichnis}
\listoftables
\cleardoublepage

\pagenumbering{arabic}
%###########################################################################
%
%   Einleitung
%
%###########################################################################
\chapter{Einleitung}
Die Aufgabe bestand darin ein Handheld zu konzipieren mit welchem es möglich ist die umgebenden \ac{GSM} \ac{BTS} zu scnannen. Durch das Auslesen des \ac{BCCH} können die gefundenen Basisstationen beschrieben und charakterisiert werden.

Ein solches System bestand bereits allerdings lief dieses unter Ubunutu und war somit an ein Notebook gebunden. Die Idee war es das Ganze ein wenig mobiler zu gestalten und auch auf den neuesten Stand zu bringen. 

Die Verwendung eines DVB-T Sticks sowie der Software Gnuradio wurde vorgeschrieben 

\chapter{Installationsanleitung}

Ausführlichere Anleitung und weitere Wiki-Einträge sind hier nachzulesen:\\
\url{https://github.com/ptrkrysik/gr-gsm/wiki/Installation-on-RaspberryPi-3}


\section{Installation von benötigter Software für gr-gsm}

\subsection{Installieren von Kalibrate}

Als erstes installieren wir Kalibrate:

\begin{code}[	numbers=left,stepnumber=1]
sudo apt-get install libtool autoconf automake libfftw3-dev librtlsdr0 librtlsdr-dev libusb-1.0-0 libusb-1.0-0-dev
git clone https://github.com/asdil12/kalibrate-rtl.git
cd kalibrate-rtl
git checkout arm_memory
./bootstrap
./configure
make
sudo make install 
\end{code}

\subsection{Zugriff auf USB-Device freischalten}

Das RTL-SDR Gerät einstecken und ID mit dem Befehl \colorbox{lightgray}{\emph{lsusb}} überprüfen. Zu sehen sollte etwas wie folgend sein:

\begin{code}
Bus 001 Device 004: ID **0bda:2832** Realtek Semiconductor Corp. RTL2832U DVB-T
Bus 001 Device 003: ID 0424:ec00 Standard Microsystems Corp. SMSC9512/9514 Fast Ethernet Adapter
Bus 001 Device 002: ID 0424:9514 Standard Microsystems Corp.
Bus 001 Device 001: ID 1d6b:0002 Linux Foundation 2.0 root hub
\end{code}

\noindent{}\\In unserem Fall ist die ID \textbf{0bda:2832}. Anschließend öffnen wir eine rules-Datei:

\begin{code}
sudo nano /etc/udev/rules.d/20.rtlsdr.rules
\end{code}

\noindent\\...in welcher dann folgende Zeile hinzugefügt werden muss:

\begin{code}
SUBSYSTEM=="usb", ATTRS{idVendor}=="0bda", ATTRS{idProduct}=="2832", GROUP="adm", MODE="0666", SYMLINK+="rtl_sdr"
\end{code}

\noindent\\Falls mehrere RTL-SDR Geräte verwendet werden, können mehrere Zeilen hinzugefügt werden. Die ID muss jeweils natürlich entsprechend des Gerätes abgewandelt werden.\\
\\Danach sollte der Raspberry Pi neugestartet werden: \colorbox{lightgray}{sudo reboot}

\subsection{Kalibrierung des RTL-SDR Gerätes}
Jetzt können wir den Befehl ausführen um das RTL-SDR Gerät zu kalibrieren (um genau zu sein um den durschnittlichen absoluten Fehler in ppm zu berechnen):

\begin{code}
kal -s GSM900 
\end{code}

\noindent\\Das Ergebnis sollte ähnlich zu diesem sein:

\begin{code}
Found 1 device(s):
  0:  Generic RTL2832U

Using device 0: Generic RTL2832U
Found Rafael Micro R820T tuner
Exact sample rate is: 270833.002142 Hz
kal: Scanning for GSM-900 base stations.
GSM-900:
	chan: 1 (935.2MHz - 33.430kHz)	power: 55085.23
	chan: 3 (935.6MHz - 34.130kHz)	power: 63242.36
	chan: 5 (936.0MHz - 33.970kHz)	power: 41270.82
...
...
	chan: 112 (957.4MHz - 32.934kHz)	power: 498930.07
	chan: 116 (958.2MHz - 31.859kHz)	power: 88039.44
	chan: 124 (959.8MHz - 32.429kHz)	power: 247404.23
\end{code}

\noindent\\Das stärkste Signal wäre in diesem Fall Kanal 112. Also führen wir die Kalibrierung auf diesem Kanal durch:

\begin{code}
kal -c 112 
\end{code}

\noindent\\...und erhalten ein Ergebnis wie folgt:

\begin{code}
Found 1 device(s):
  0:  Generic RTL2832U

Using device 0: Generic RTL2832U
Found Rafael Micro R820T tuner
Exact sample rate is: 270833.002142 Hz
kal: Calculating clock frequency offset.
Using GSM-900 channel 112 (957.4MHz)
average		[min, max]	(range, stddev)
- 34.368kHz		[-34376, -34357]	(20, 4.697051)
overruns: 0
not found: 0
average absolute error: 35.897 ppm
\end{code}

\noindent\\Unser durschschnittlicher absoluter Fehler wäre hier also 36 ppm (35.897 ppm).

\subsection{Installation von GNU Radio}

Als nächstes installieren wir GNU Radio:

\begin{code}
sudo apt-get install gnuradio gnuradio-dev
\end{code}

\subsection{Installieren von libosmocore}

Libosmocore muss kompiliert werden...

\begin{code}[numbers=left, stepnumber=1]
sudo apt-get install cmake
sudo apt-get install build-essential libtool shtool autoconf automake git-core pkg-config make gcc
sudo apt-get install libpcsclite-dev libtalloc-dev
git clone git://git.osmocom.org/libosmocore.git
cd libosmocore/
autoreconf -i
./configure
make
sudo make install
sudo ldconfig -i
cd
\end{code}

\noindent\\...außerdem benötigen wir noch ein paar andere Dinge.

\begin{code}[numbers=left, stepnumber=1]
sudo apt-get install gr-osmosdr rtl-sdr
sudo apt-get install libboost-dev
sudo apt-get install osmo-sdr libosmosdr-dev
sudo apt-get install libusb-1.0.0 libusb-dev
sudo apt-get install libboost-all-dev libcppunit-dev swig doxygen liblog4cpp5-dev python-scipys
\end{code}

\section{Installation von gr-gsm}

Und nun zum letzten Schritt:

\begin{code}[numbers=left, stepnumber=1]
git clone https://github.com/ptrkrysik/gr-gsm.git
cd gr-gsm
mkdir build
cd build
cmake ..
make
sudo make install
sudo ldconfig
\end{code}

\noindent\\Zuletzt erstellen wir noch die ~/.gnuradio/config.conf config-Datei, mit \colorbox{lightgray}{\emph{nano ~/.gnuradio/config.conf}}. Und fügen diese zwei Zeilen hinzu (damit GNU Radio die custom Blöcke von gr-gsm finden kann):

\begin{code}
[grc]
local_blocks_path=/usr/local/share/gnuradio/grc/blocks
\end{code}

\chapter{Anleitung zur GUI unter PyQt4}

Im folgenden Kapitel soll die GUI und deren Inhalt anhand einiger Bildbeispiele verdeutlich werden.
Nachfolgend wird erläutert wie und wo im Code es möglich ist, diese Inhalte anzupassen bzw. zu verändern.

\section{Allgemeiner Aufbau der GUI}

Starten des grgsm-scanner initiiert die GUI wie zu sehen in Abbildung:\ref{fig1}.
Das GUI-Fenster ist in zwei Felder aufgeteilt, ein kleineres Feld auf der linken Seite
und ein größeres auf der Rechten.
Im linken Feld werden die Kanäle dargestellt, welche nach und nach gescannt und hinzugefügt werden.
Im rechten Feld wird zu Beginn Informationen zur Projektarbeit ausgegeben. Diese wird durch die Informationen von ausgewählten Kanälen überschrieben, dazu jedoch später mehr.

\mydennis{GUI_Window}{Standard Fenster GUI}{fig1}{0.0}{0.9\textwidth}

\noindent Solange die Suche nach Kanälen noch nicht abgeschlossen wurde, wird oben im Fensterrahmen \textbf{-searching for GSM BTS-} dargestellt (siehe Abbildung:\ref{fig1}).
Falls der Scan abgeschlossen ist, wird dies durch ein \textbf{Done} angezeigt, ersichtlich aus Bild \ref{fig2}. Außerdem sollten nun alle gefundenen Kanäle im linken Feld zu sehen sein.

\mydennis{GUI_Done}{Abgeschlossener Scan}{fig2}{0.0}{0.9\textwidth}


\noindent Falls Kanäle schon gescannt wurden und somit links auftauchen, können diese ausgewählt werden um Informationen der einzelnen Parameter einzusehen.
Unter den Parametern befinden sich der ARFCN, Frequenz, LAC, MCC, MNC, CID und die Power.
Zusätzlich werden ein paar weitere Informationen gegeben (siehe Abbildung:\ref{fig3}).
Sollte ein Maus angeschlossen sein können die Informationen auch durch eine Art Tooltip über mouse-hovering auf dem jeweiligen Kanal ausgegeben werden (siehe Abbildung:\ref{fig4}).

\mydennis{GUI_Textinfo}{Info des BCCH}{fig3}{0.0}{0.9\textwidth}

\mydennis{GUI_Tooltip}{Mouseover Tooltip}{fig4}{0.0}{0.9\textwidth}


\newpage Bei jedem Kanal ist es möglich durch einen kleinen Pfeil rechts daneben, diesen aufzurollen um die Parameter einzeln zu betrachten (siehe Abbildung:\ref{fig5}). Durch anklicken wieder auf die einzelnen Parameter wird die jeweilige Information dargestellt (siehe Abbildung:\ref{fig6}).

\mydennis{GUI_aufgeklappt}{GUI Ansicht angeklickt}{fig5}{0.0}{0.9\textwidth}

\mydennis{GUI_childinfo}{Infos der Childnodes}{fig6}{0.0}{0.9\textwidth}


\section{Anleitung zum Anpassen des textuellen Inhalts}

\subsection*{Anpassung der Titelseite}
Die Titelseite, wie im vorherigen Kapitel bereits erwähnt, welche zu sehen ist am Programmstart, kann durch Änderungen im folgenden Programmcode (Listing: \ref{code1}) angepasst werden.
Die Zeilenangaben links verdeutlichen zusätzlich, wo man dies im Programmcode finden kann.

\begin{code}[firstnumber=100,numbers=left,stepnumber=1, caption={Titelseite},captionpos=b,label={code1}]
self.text = QtGui.QLabel("\n"
                         " \t\t   GrGsm Scanner GUI\n\n\n"
                         " Mit Hilfe der GNURADIO Bloecke \"GrGsm\" von Piotr Krysik\n"
                         " wird der BCCH umliegender GSM BTS ausgelesen und dargestellt\n\n"
                         " Ein Klick auf den gewuenschten Kanal links zeigt die Parameter an\n\n\n\n\n\n\n\n\n\n\n"
                         " Projektarbeit unter Prof. Litzenburger im SoSe 2017\n"
                         " von Dennis Dette und Christian Kobiela\n\n"
                         , self.right)
\end{code} 

\subsection*{Anpassung der Parentnode}

Die Informationen, welche angezeigt werden wenn man einen gefunden Kanal anklickt, können in den folgenden Codezeilen geändert werden (Listing: \ref{code2}). Bisher werden wie zu sehen ist die einzelnen Parameter mit Werten dargestellt und Angaben über Land und Betreiber der BTS. 

\begin{code}[firstnumber=187,numbers=left,stepnumber=1, caption={Kanalinformation},captionpos=b,label={code2}]
parent_item.setToolTip(column,"Der Kanal mit der " + Buffer_Channel.ARFCN + 
		" enthaelt folgende Informationen: \n\n"
		+ " " +Buffer_Channel.ARFCN + " \t\tAbsolute Radio Frequency Chan. No.\n"
		+ " " +Buffer_Channel.FREQ + " \t\tFrequency in Mhz\n"
		+ " " +Buffer_Channel.LAC + " \t\tLocation Area Code\n"
		+ " " +Buffer_Channel.MCC + " \t\tMobile Country Code\n" 
		+ " " +Buffer_Channel.MNC + " \t\tMobile Network Code\n"
		+ " " +Buffer_Channel.CID + " \t\tCell Identification\n"			                        
		+ " " +Buffer_Channel.PWR+ " \t\tPower\n\n\n"
		+ " Laut " + Buffer_Channel.MNC + " ist der Betreiber " + Betreiber
		+ "\n Laut " + Buffer_Channel.MCC + " steht die BTS in " + str(Land)
		)
\end{code}

\subsection*{Anpassung der  Childnodes}

Nun zum letzten Teil, dessen Text angepasst werden kann, den Parameterinformationen. Diese sind einzusehen, wenn ein Kanal aufklappt ist. Hier werden durch auswählen der einzelnen Parameter die Informationen dargestellt, welche im Code \ref{code3} zu sehen sind. Die sieben Parameterinformationen werden in den verschiedenen If-Fällen definiert. 

\begin{code}[firstnumber=204,numbers=left,stepnumber=1,caption={Parameterinformationen},captionpos=b,label={code3}]
if k==0:
	item.setToolTip(column, " Absolute Radio Frequency Channel Number:\n\n " + Buffer_Channel.ARFCN 
	+ "\n\n In GSM cellular networks, an absolute radio-frequency channel number\n ""(ARFCN) is a code that specifies a pair of physical radio carriers used for\n ""transmission and reception in a land mobile radio system, one for the uplink\n ""signal and one for the downlink signal.\n\n "
	"ARFCN for GSM 900 \n "
	"ARFCN = (Chan. Freq - 45 Mhz- 890 Mhz)/200\n\n "
	"ARFCN for E GSM 900\n "
	"ARFCN = 1024 + (Chan. Freq - 45 Mhz - 890 Mhz)/200\n ")
elif k==1:
	item.setToolTip(column, " Frequency:\n\n " + Buffer_Channel.FREQ + "\n\n GSM Downlink Frequency")
elif k==2:
	item.setToolTip(column, " Cell Identification:\n\n " + Buffer_Channel.CID +
	"\n\n A GSM Cell ID (CID) is a generally unique number used\n "
	"to identify each base transceiver station (BTS) or\n "
	"sector of a BTS within a location area code (LAC) \n "
	"if not within a GSM network.")				
elif k==3:
	item.setToolTip(column, " Location Area Code:\n\n " + Buffer_Channel.LAC+
	"\n\n A location area is a set of base stations that are\n "
	"grouped together to optimise signalling.\n "
	"To each location area, a unique number called a location area code is assigned.\n ")
				    
					
elif k==4:
	item.setToolTip(column, " Mobile Country Code:\n\n " + Buffer_Channel.MCC +
	"\n\n The mobile country code consists of 3 decimal digits and the mobile\n "
	"network code consists of 2 or 3 decimal digits. The first \"2\" in 262\n "
	"stands for Europe, 62 for Germany")
				
elif k==5:
	item.setToolTip(column, " Mobile Network Code:\n\n" + Buffer_Channel.MNC +
	"\n\n Legende:\n 01: Telekom Deutschland\n 02: Vodafone\n 03: Telefonica (ehem. E-Plus)\n 10: DB Telematic")
	
elif k==6:
	item.setToolTip(column, " Power:\n\n " + Buffer_Channel.PWR + 
	"\n\n GSM needs at least a Power of -102dBm " )
\end{code}
 
%###########################################################################
%
%   Kapitel 1
%
%###########################################################################

\chapter{Bedienung des Handhelds}
In diesem Kapitel finden Sie eine Bedienungsanleitung für das Handheld und Tipps für den Umgang mit diesem.

Um einen gemeinsamen Betrachtungspunkt zu haben legen wir nun fest, dass man das Tablet im Querformat verwendet und "oben" die Seite beschreibt an der die USB Anschlüsse des Raspberrys zu sehen sind. In diesem Fall finden wir die Ladebuchse rechts über die das Tablet sowohl geladen als auch stationär verwendet werden kann. HDMI, AUX und einen Mikro USB Anschluss 2 zur Direktversorung, Überbrückung des Akkus, links. Der ON/OFF Schalter befindet sich auf der Unterseite rechts. Dieser ist leider nicht all zu gut zugänglich, weshalb es empfohlen wird den Schalter mit einem spitzen Gegenstand zu betätigen (Bspw. bietet sich hier die Antenne an).

\mywidthpicture{Draufsicht}{Beschreibung der Anschlüsse}{fig:Draufsicht-fig}{0.0}{0.8\textwidth}

\section{Stromversorgung}

Verbaut ist ein 2,5 Ah LiPo Akku welcher über eine Adafruit Powerboost1000c Ladeelektronik geladen und betrieben wird. Die Elektronik ist sowohl dafür zuständig den Akku aufzuladen als auch bei der stationären Verwendung ein angeschlssossenes USB Netzteil als Stromquelle zu Nutzen. Da der Normalstrom der aus dem Akku gezogen wird sich um die 1A bewegt kann es durchaus zu einer Unterversorgung kommen. Aufgrund dessen wird es nicht empfohlen bei höhster Displayhelligkeit den GSM Suchlauf durchzuführen. Die Displaybeleuchtung dunkelt sich nach 10 Sekunden ab um diesen Fall der möglichen Unterversorgung auszuschließen. Dies kann in den Einstellungen geändert werden. 

\mywidthpicture{Displayhelligkeit}{Einstellung der automatischen Helligkeitsregelung}{fig:Dislpayhelligkeit-fig}{0.0}{\textwidth}

\section{Bedienung über die Desktopoberfläche}

Auf dem Desktop befinden sich die wichtigsten Shortcuts für den Gebrauch des GSM Scanners. 

\mywidthpicture{Desktop}{Ansicht des Dekstops}{fig:Desktop-fig}{0.0}{\textwidth}


Das Touchscreen ist so eingestellt, dass man nur einmal klicken muss um Programme auszuführen was eine Bedienung mit dem Finger erleichtern sollte. Um eine einfache Einstellbarkeit der Displayhelligkeit zu realisieren haben wir auf dem Desktop Shortcuts hierfür implementiert. Die Stufen 10\%, 50\% und 100\% können ausgewählt werden. Sind andere Stufen gewünscht, so kann man die Displayhelligkeit durch ausführen des Befehls
\begin{code}
echo XXX > /sys/class/backlight/rpi_backlight/brightness
\end{code}
\noindent\\ändern. XXX kann im Bereich von 0 (0\%) bis 255 (100\%) gewählt werden. 

Ferner findet sich ein virtuelles Keyboard auf dem Desktop falls man mobil etwas schreiben möchte. Reboot und Shutdown Shortcuts sind ebenso zu finden. Bitte beachten Sie: Nach dem Shutdown muss die Stromversorgung zusätzlich am ON/OFF Schalter getrennt werden.

Der GSM Scanner hat ebenfalls ein Desktop Shortcut welches mit einem Klick die Suche nach GSM Basisstationen ermöglicht. Um den Hintergrund zu verstehen erläutere ich hier auf die möglichen Einstellmöglichkeiten mit denen ein Scan gestartet werden kann. 

\section{Erweitern und Verändern}

Nach dem Anschalten am ON/OFF Schalter auf der Unterseite des Handhelds fährt dieses hoch ohne, dass eine Anmeldung erforderlich ist. Sollte dies geändert werden, so müssen müssen zwei Befehle in der Datei lightdm.conf auskommentiert werden. 
\begin{code}
cd /etc/lightdm
nano lightdm.conf

****************
autologin-user=root
autologin-user-timeout=0
****************

^ muessen durch voransetzen eines "#" auskommentiert werden
\end{code}


\begin{code}
Options:
  -h, --help            show this help message and exit
  -b BAND, --band=BAND  Specify the GSM band for the frequency. Available
                        bands are: GSM900, DCS1800, GSM850, PCS1900, GSM450,
                        GSM480, GSM-R
  -s SAMP_RATE, --samp-rate=SAMP_RATE
                        Set sample rate [default=2000000.0] - allowed values
                        even_number*0.2e6
  -p PPM, --ppm=PPM     Set frequency correction in ppm [default=0]
  -g GAIN, --gain=GAIN  Set gain [default=24.0]
  --args=ARGS           Set device arguments [default=]
  --speed=SPEED         Scan speed [default=4]. Value range 0-5.
  -v, --verbose         If set, verbose information output is printed: ccch
                        configuration, cell ARFCN's, neighbour ARFCN's
                        
                   
grgsm_scanner -g 50 -p 29 --speed=5
\end{code}
\noindent\\Mit diesen Übergabeparametern wird der GSM Scanner aufgerufen und sucht im default BAND GSM900 das Netz von 925 Mhz bis 960 Mhz ab. Dies entspricht dem E-GSM 900 Netz. Ein Gain von 50 dB ist der Maximalwert und ermöglicht somit die größte "Ausbeute" was die Ergebnisliste angeht. Der Offset des Quarzes bei Betriebstemperatur wurde berechnet durch eine Kalibierung des DVB-T Sticks und ist immer mit anzugeben. 

Das zugehörige Desktop Shortcut, wie auch die Shortcuts zur Helligkeitsregulierung, führen Shell Skripte aus welche unter 
\begin{code}
/root/Documents
\end{code}
\noindent\\hinterlegt sind. Hier kann man eingreifen falls etwas geändert werden soll. Am besten öffnet man diese über das Terminal mit "nano" da sonst kein Textverarbeitungsprogramm installiert ist. 

Alle von uns geschriebenen Komponenten sind im Ordner 
\begin{code}
/root/GrGsm-Gui
\end{code}
\noindent\\zu finden. Da es sich hierbei um ein Git Repository handelt kann dieses auch über 
\begin{code}
cd GrGsm-Gui/
git pull
\end{code}
\noindent\\auf den neusten Stand gebracht werden, sollten Veränderungen vorgenommen werden. 

%###########################################################################
%
%   Kapitel 3
%
%###########################################################################

\chapter{Hardwareaufbau des Handhelds}

Im folgenden wird erklärt wie die verbauten Komponenten zusammenspielen. Die Anleitung für das Case stammt von Adafruit.com und wird als PDF Datei der DVD beiliegen. 

\section{Stromversorgung}

\mywidthpicture{Verkabelung}{Verkabelung im Inneren des Case}{fig:Verkabelung-fig}{0.0}{\textwidth}

\section{Benötigten Teile}
\subsection{Hardware}
\begin{itemize}
\item Pi Foundation PiTFT - 7" Touchscreen Display 
\item Raspberry Pi 3 
\item 200mm Flex Displaykabel
\item Adafruit PowerBoost 1000C
\item 2500mAh LiPo Akku
\item SPDT Schalter
\end{itemize} 

\subsection{Werkzeuge und Ergänzendes}
\begin{itemize}
\item 3D Drucker 
\item Filament  
\item Kreuzschlitzschraubenzieher 
\item Lot 
\item Litzen mit 1,5 mm$^2$
\item M3 x .5 x 6M Schrauben 
\item 2-56 3/8 Schrauben 
\end{itemize}

%###########################################################################
%
%   Kapitel 2
%
%###########################################################################
\chapter{Setup des Raspberry PI}
Im folgenden wird beschrieben wie wir den Raspberry aufgesetzt haben und welche Probleme uns dabei begegnet sind.
\section{Das How To und interdependencies}

GrGsm wird hauptsächlich von Piotr Krysik entwickelt und wird von diesem auch gepflegt. Es gibt auch einen vorgeschlagenen Weg wie man GrGsm auf einem Raspberry Pi 3 zu installieren hat. Dieser Weg hat sich allerdings als Sackgasse erwiesen. Grund dafür war, dass in Zwischenzeit GnuRadio weiterentwickelt wurde und GrGsm sich dem angepasst hatte. Die Weiterentwicklung betraf leider nicht Raspbian Jessie was zu folgenden Problemen führte.

GrGsm benötigte damals eine Version von GnuRadio > 3.7.9. Für Raspbian stand allerdings nur die Version 3.7.6 zur Verfügung. Der Versuch auf instabile Testversionen 3.7.10+ auszuweichen erwies sich ebenso als aussichtslos. Das Installieren von diesen Versionen führte nur dazu, dass wiederum Bibliotheken von denen GnuRadio abhängig war inkompatibel wurden. 

Nach einiger Recherche wie man ein funktionierendes Gesamtpaket erhalten könnte sind wir auf das Programm PyBombs gestoßen. Dieses macht genau das, es installiert ein Programm und alle Abhängigkeiten die dieses zur Nutzung benötigt. Im Gegensatz zu apt-get, wie man es von von Ubuntu und sonstigen Linux Distributionen kennt bezieht sich PyBOMBS nicht nur auf die Standard Repositorys und die dort hinterlegte Softwareversionen sondern installiert genau die Versionen die benötigt werden. Hierbei kann sich PyBOMBS auf eigens hinterlegte Repositorys, sogenannte Recipes beziehen von wo es die Software aus den Quelldateien kompiliert. 
So schön es klingt führte es zu den selben Problemen wie im ersten Versuch nur, dass dieses mal andere Programme inkompatibel wurden.


An diesem Punkt haben wir beschlossen, dass eine zeitnahe Umsetzung mit GrGsm wohl aussichtslos erscheint. Gleichzeitig standen wir mit Piotr Krysik in Kontakt um einen Weg zu finden wie es doch zu lösen sein könnte. Seither ist Herr Krysik dabei seine Anpassungen in einem Raspberry Emulator zu überprüfen. Wir entschieden uns dennoch Airprobe zu versuchen, da es nicht absehbar war wann wir mit einer Antwort hätten rechnen können. Dafür haben wir uns an einer jahrealten Anleitung orientiert und sind dabei auf Kali Linux gestoßen, dieses war glücklicherweise auch für Raspberrys verfügbar. Nachdem wir das installiert hatten, war, wer hätte es gedacht Gnuradio mittlerweile zu Neu um mit Airprobe arbeiten zu können. Allerdings war die stabile Version glücklicherweise 3.7.10, somit mit dem neusten GrGsm kompatibel und so konnten wir alle Abhängigkeiten installieren, die Kalibrierung durchführen und GrGsm installieren. Zu der Zeit haben wir schon einen Monat Arbeit in das Projekt stecken und die Grenzen unserer Frustrationstoleranz immer ein Stückchen weiter schieben müssen.

Im Endeffekt ist die Anleitung wie sie hier in Kapitel \ref{Install} dokumentiert wird, unter Kali Linux voll durchführbar. Es wäre allerdings wünschenswert den GSM Scanner unter Raspbian zum laufen zu bringen da es sich hierbei um ein besser auf den Raspberry zugeschnittenes Betriebssystem handelt. 
Damit diese Portierung auf Raspbian möglich ist, müssen die vorher beschriebenen Probleme mit den Abhängigkeiten gelöst werden. 

%###########################################################################
%
%   Kapitel 1
%
%###########################################################################
\chapter{Einige Features}
Hier werden einige Features gezeigt.

\section{Ein Unterabschnitt}
Die Tabelle~\ref{tab:example-table} ist ein Beispiel für eine wissenschaftliche Tabelle (ohne vertikale Trennlinien).

\begin{table}[ht]
\caption{Beispieltabelle}
\centering
\begin{tabular}{@{}ccc@{}} \toprule
    Apfel & Birne & Katze \\
    \SI{200}{g} & \SI{180}{g} & \SI{3.4}{kg} \\ \bottomrule
\end{tabular}
\label{tab:example-table}
\end{table}

So zitiert man \cite{slam06}. Man kann auch mehrere Quellen auf einmal referieren \cite{pomax16, dlib09, biagiotti_open-loop_2012, hunter_matplotlib:_2007}

In Abbildung~\ref{fig:example-fig} ist eine Kurve dargestellt.
\mywidthpicture{fail_path_linearization}{Beispielbild}{fig:example-fig}{0.0}{\textwidth}

Inkscape erlaubt es eps-Files und den Text getrennt zu exportieren, auf diese Weise kann man die Schrift nachträglich anpassen. Ein Beispiel dazu ist in Abbildung~\ref{fig:inkscape-fig} dargestellt.
\myinkscapepicture{mr_kinmatic_model.eps_tex}{Inkskape Bild mit Text}{fig:inkscape-fig}{0.6\textwidth}

Außerdem gibt es eine Umgebung zum Schreiben von Pseudo-Code. Ein Beispiel dazu ist in Algorithmus~\ref{alg:example-alg} dargestellt.
\begin{algorithm}
\begin{algorithmic}[1]
\Procedure{AddTwoNumbers}{$x$, $y$}
	\State $sum \gets x + y$
   	\If{$sum = 42$}
		\State \Call{FixEverything}{}
   	\EndIf
\EndProcedure
\end{algorithmic}
\caption{Beispielalgorithmus}
\label{alg:example-alg}
\end{algorithm}

Das Paket acronym handelt Abkürzungen automatisch. Bei der ersten Verwendung sieht das so aus: \ac{UML}. Wenn die Abkürzung nochmal verwendet wird, steht da nur noch \ac{UML}.



%###########################################################################
%
%   Zusammenfassung und Ausblick
%
%###########################################################################
\chapter{Zusammenfassung}

In dieser Projektarbeit ist es uns gelungen ein Handheld zum Auslesen des BCCH zu konzipieren, diesen aufzubauen und in Betrieb zu nehmen. Als Herzstück dient uns ein von Kali Linux betriebener Raspberry Pi 3, das dazugehörige Display mit Toucheingabe und ein DVB-T Stick. Letzterer wurde nach dem Prinzip des Software Defined Radio verwendet, welches besagt, dass möglichst viel der Signalverarbeitung in der Software durchgeführt wird. So können die analogen Komponenten der Kommunikationssysteme einfach gehalten werden.

Mit den Gnuradio Blöcken "GrGsm" von Piotr Krysik ist es uns gelungen den Informationsstring des BCCH auszulesen und mit einer eigens geschriebenen GUI darzustellen. Die Darstellung der GUI und der Suchvorgang der GSM Basisstationen laufen parallel auf zwei Threads um eine laufende Aktualisierung zu ermöglichen.

Mit dem fertigen Handheld ist eine mobile Nutzung von bis zu 5 Stunden möglich. So kann auch die Netzabdeckung der entferntesten Orte analysiert werden.
%###########################################################################
%
%   Ausblick
%
%###########################################################################
\chapter{Ausblick}

\addcontentsline{toc}{chapter}{Literaturverzeichnis}
\bibliographystyle{unsrtdin}
\bibliography{my_literature}


\end{document}

%##########################################################################
